% from ega3-1-fr.tex
\oldpage[III]{81}
\phantomsection
\label{chapter:III-fr}

\begin{center}
{\large CHAPITRE III\par}
\vspace{1.5em}
{\Large\bfseries ÉTUDE COHOMOLOGIQUE\par}
{\Large\bfseries DES FAISCEAUX COHÉRENTS\par}
\vspace{1.8em}
{\bfseries Sommaire\par}
\end{center}

\medskip
\begin{tabular}{@{}r@{\ }p{0.82\textwidth}@{}}
\S~1. & Cohomologie des schémas affines.\\
\S~2. & Étude cohomologique des morphismes projectifs.\\
\S~3. & Le théorème de finitude pour les morphismes propres.\\
\S~4. & Le théorème fondamental des morphismes propres~; applications.\\
\S~5. & Un théorème d'existence de faisceaux algébriques cohérents.\\
\S~6. & Foncteurs Tor et Ext locaux et globaux~; formule de Künneth.\\
\S~7. & Étude du changement de base dans les foncteurs cohomologiques
covariants de Modules.\\
\S~8. & Le théorème de dualité sur les fibrés projectifs.\\
\S~9. & Cohomologie relative et cohomologie locale~; dualité locale.\\
\S~10. & Relations entre cohomologie projective et cohomologie locale.
Technique de complétion formelle le long d'un diviseur.\\
\S~11. & Groupes de Picard globaux et locaux\footnotemark.
\end{tabular}

\medskip
Ce chapitre donne les théorèmes fondamentaux sur la cohomologie des faisceaux
algébriques cohérents, à l'exception de ceux découlant de la théorie des
résidus (théorèmes de dualité), qui feront l'objet d'un chapitre ultérieur.
Parmi les premiers, il y a essentiellement six théorèmes fondamentaux,
faisant l'objet des six premiers paragraphes du présent chapitre. Ces
résultats seront dans la suite des outils essentiels, même dans des questions
qui ne sont pas de nature proprement cohomologique~; le lecteur en verra les
premiers exemples dès le \S~4. Le \S~7 donne des résultats de nature plus
technique, mais d'un usage constant dans les applications. Enfin, dans les
\S\S~8 à 11, nous développerons certains résultats, liés à la dualité des
faisceaux cohérents, particulièrement importants pour les applications, et
qui peuvent s'exposer antérieurement à la théorie générale des résidus.

\footnotetext{Le chapitre IV ne dépend pas des \S\S~8 à 11, et sera sans doute
publié avant ces derniers.}

